\subsection{Uniform dyadic dominance and the strict cubic indices}

\begin{theorem}[P-F57: the selected dyadic pair dominates]
\label{thm:dyadicdominance}
For the ordinary cubic plethysm family $f_n$ of P-F31, the
quantities in (P-F56l) satisfy
\[
                         \frac{E_n}{|Z_n|}\longrightarrow0.
 \tag{P-F57a}
\]
Consequently, for every fixed real $0<c<1/4$ and all
sufficiently large $n$,
\[
                  p(f_n)>2^{cn},
             \qquad s(f_n)\ge2^{cn}.
\]
These are the literal positivity and strict-saturation indices
of GCT VI, Definitions~1.2.2 and 1.2.4.
Since the source parameter is
$h=\binom{n+2}{3}^{\!2}+1+n=\Theta(n^6)$, both narrow
polynomial-index hypotheses fail on this ordinary irreducible
cubic-inner family.
\end{theorem}

\begin{proof}
We prove the branch contraction (P-F56ad), rather than
extrapolating from finite computations.
Use $M=2^{n-1}$, $l=n-1$, the natural selected residue
$b_n=3^{-1}\pmod M$, and the positive increment product
$J_n(a)$ from (P-F56ab)--(P-F56ac).
The comparison has a near part, proved next, and a far part
proved in the layer-cake lemma below.

\paragraph{Near increments.}
For odd $1\le j<M$, let
$a_j\equiv j3^{-1}\pmod M$ and
$R=\lceil\log_2(16j)\rceil$. For all $n\ge5$,
\[
 \left|
 \log_2\frac{J_n(a_j)}{J_n(b_n)}
 +\binom{n-1}2\log_2j
 -\mathbf1_{3\mid j}\binom n2
 \right|\le100nR^2.
 \tag{P-F57b}
\]
We supply the product estimates, including the error uniform
in the multiplier.

The increment multiset consists of the bulk weights
$1+w_u+w_v$ with $2\le u<v\le l-1$, except $(u,v)=(2,3)$;
the weights $2+w_k$ with $3\le k\le l-1$; and the weight
six, all divisible by three. Their total number is
\[
 B=\left(\binom{n-3}2-1\right)+(n-4)+1
                         =\frac{(n-1)(n-4)}2.
 \]
The exceptional weights are
$1$, $2$, and $1+w_k$ for $2\le k\le l-1$.
There are $n-1$ of them.
Give each weight its largest active dyadic index $k$ and set
$q=l-k$. Thus the weights one and two have $q=l$ and
$q=l-1$, respectively.
The exceptional weights have exactly one occurrence of
each $q=1,\ldots,l$, so
\[
 B+(n-1)=\binom{n-1}2,\qquad
 \sum_{\rm exceptional}q=\frac{l(l+1)}2=\binom n2.
 \tag{P-F57c}
\]
For a divisible-by-three weight, the absolute value of the
integer representative lies between $2^{k-2}$ and
$2^{k+1}$. This is immediate for six and $2+w_k$.
For $1+w_u+w_v$, it follows from
$2^v-2^{v-1}-1\ge2^{v-2}$, since $v\ge3$.
For each fixed $q$, there are at most $2n$ such weights.
An exceptional weight has absolute value between
$2^{k-1}$ and $3\cdot2^{k-1}$, with the same bounds
also valid for the singleton weights one and two.

Write
$\eta_j=\exp(2\pi ia_j/M)
        =\omega_j\exp(2\pi ij/(3M))$.
As in (P-F56ap), $\omega_j^3=1$ and
$\omega_j=1$ exactly when $3\mid j$.
Call a factor low if $q>R$ and high otherwise.
The sparse dyadic separation lemma (P-F56an) applies
to every factor: each weight is a signed sum of at most
three distinct powers of two.
In particular a high target factor is at least
$2^{-3R-1}$, and at least $2^{-2R-1}$ for an
exceptional factor, which uses at most two powers.
All factors are at most two.

For a low divisible-by-three factor, there is no cube-root
shift. Its target phase is $j$ times its reference phase,
and its absolute target phase is at most
$(2j/3)2^{-q}<1/24$.
The sine estimate proved in the active-triple part of P-F56
shows that the logarithm of its factor ratio differs from
$\log_2j$ by at most $16j^2 4^{-q}$.
Summing over at most $2n$ weights at each $q>R$ gives
error less than $n$, since
$j^2 4^{-R}\le1/256$.
For a high divisible-by-three factor, the reference phase
has absolute value between $2^{-q}/12$ and
$2\cdot2^{-q}/3<1/2$.
Its factor is at least $2^{-q}/3$.
The high target lower bound above therefore bounds the
absolute logarithm of the factor ratio by $5R$.
There are at most $2nR$ high factors.
Restoring their missing $\log_2j$ main terms costs
at most $2nR^2$.
Thus the contribution of all divisible-by-three factors
to the logarithm of the reciprocal product is
$-B\log_2j$ with error at most $20nR^2$.

Every exceptional reference factor lies between $1/3$ and
two. Indeed its cube-root shift is nontrivial, and the
perturbation of the phase has absolute value at most $1/4$.
If $3\nmid j$, every low target exceptional factor has
the same nontrivial type of shift, perturbed by less than
$1/24$, so its logarithmic ratio has absolute value
at most three.
There are at most $R$ high exceptional factors.
Their two-term sparse separation bound makes each
logarithmic ratio at most $3R$ in absolute value.
The total exceptional contribution is then bounded by
$10nR^2$ in absolute value.

If $3\mid j$, the target shift instead vanishes.
A low exceptional target factor lies between
$(2/3)j2^{-q}$ and $4j2^{-q}$.
Its reciprocal-factor ratio therefore has logarithm
$q-\log_2j$ with error at most four.
Restore the $q$ main terms of the at most $R$ high
factors: their sum is at most $R(R+1)/2\le R^2$.
The actual high logarithmic ratios total at most $3R^2$.
Discarding all the $-\log_2j$ terms costs at most $nR$,
and the low constant errors cost at most $4n$.
Together these bounds show that the exceptional contribution
is $\binom n2$ with error at most $10nR^2$.
We have proved, in both cases,
\[
 \log_2\frac{J_n(a_j)}{J_n(b_n)}
 =-B\log_2j+\mathbf1_{3\mid j}\binom n2+\mathcal E,
 \qquad |\mathcal E|\le30nR^2.
 \tag{P-F57d}
\]
Replacing $B$ by $\binom{n-1}2$ costs only
$(n-1)\log_2j\le nR$.
This proves (P-F57b), with room in the stated constant.

Fix
$\delta=1/100000$.
For every odd $3\le j\le2^{\delta n}$, (P-F57b)
has the following explicit consequence:
\[
 \log_2\frac{J_n(a_j)}{J_n(b_n)}
                         \le-\frac{n^2\log_2j}{12}.
 \tag{P-F57e}
\]
To check it, put $t=\log_2j>3/2$, so $R\le t+5$.
Use
$\binom{n-1}2\ge n^2/2-2n$ and
$\binom n2\le n^2/2$ in (P-F57b).
The result is at most
\[
 -\frac{(t-1)n^2}{2}
                +n\bigl(2t+100(t+5)^2\bigr)
 \le-\frac{tn^2}{6}+2502nt^2.
 \]
Here $t-1\ge t/3$, $t+5\le5t$, and $t>1$.
Since $t\le n/100000$ and
$2502/100000<1/12$, this proves (P-F57e).
An empty near range requires no estimate.

\paragraph{P-F57, far-increment estimate.}
Use the exact increment $J_n(a)$ of (P-F56ab)--(P-F56ac),
with $n\ge5$, $l=n-1$, $M=2^l$, and $b_n=3^{-1}\pmod M$.
For an odd residue $a\pmod M$, let
\[
 J=\min\{j,M-j\},\qquad j\equiv3a\pmod M,
                 \quad 1\le j<M.
\]
Thus $J$ is an odd integer with $1\le J<M/2$.
To distinguish the integer $J$ from the increment, write
$I_n(a)=J_n(a)$ in this paragraph. Then
\[
 \log_2\frac{I_n(a)}{I_n(b_n)}
       \le 7n^2-\frac{(\log_2J-8)_+^3}{24}.
 \tag{P-F57f1}
\]
In particular, for each fixed $\delta>0$, uniformly for
$J\ge2^{\delta n}$ and all sufficiently large $n$,
\[
       \frac{I_n(a)}{I_n(b_n)}
                    \le2^{-\delta^3n^3/384}.
 \tag{P-F57f2}
\]
For example it suffices here that
$n\ge\max\{5,16/\delta,2688/\delta^3\}$.

\emph{Proof.}
Put $w_i=(-2)^i$ for $0\le i<l$ and $w_l=0$.
We first bound all the active pair factors
\[
       1-\exp\bigl(2\pi i a(1+w_j+w_k)/M\bigr),
                 \qquad 2\le j<k<l.
 \tag{P-F57f3}
\]
The increment includes all of these except $(j,k)=(2,3)$.
Every phase is nonzero modulo one, because its numerator
is odd. Write $\theta=a/M$ and
$y_{jk}=\theta(1+w_j+w_k)$ modulo one.
For an integer $u\in\{0,\ldots,l-1\}$ put
$r=l-u$ and let $N_u$ be the number of pairs in (P-F57f3)
with $\|y_{jk}\|<2^{-u}$.

The map from subsets of $\{0,\ldots,l-1\}$ to residues
modulo $2^l$ that sends a subset to the sum of its $w_i$
is injective. Indeed the difference of two distinct subset
sums has a unique least index $i$ with nonzero coefficient
$1$ or $-1$. Its two-adic valuation is exactly $i<l$.
This also proves injectivity for the pair sums used below;
the zero-weight variable $w_l$ is not included here.

Call an active index early if $j<r$, and late otherwise.
The number of early indices is $e=(r-2)_+$.
There are at most $\binom e2$ early--early pairs.
For each fixed early index $j$, all mixed pair numerators
$a(1+w_j+w_k)$ have the same residue modulo $2^r$,
since $k\ge r$ and $a$ is odd. Their residues modulo $M$
are distinct. In centered representatives, the condition
$\|y_{jk}\|<2^{-u}$ places them in the interval
$(-2^r,2^r)$, possibly intersected with the centered
interval of length $M$. Each residue class modulo $2^r$
has at most three representatives in that interval.
Thus there are at most $3e$ good mixed pairs.
All late--late pair numerators are congruent to $a$
modulo $2^r$, and distinct by the subset-sum injectivity.
There are at most three good late--late pairs. Consequently
\[
             N_u\le\binom r2+3r+3.
 \tag{P-F57f4}
\]

There is a loss when the anchor is separated from a cube
root. Set $\tau=\|3\theta\|=J/M$ and suppose
$\tau>3\cdot2^{-u}$.
Partition the $e$ early indices into
$h=\lfloor e/2\rfloor$ disjoint consecutive pairs,
starting with $(2,3)$, and discard at most one final index.
For each two distinct index-pairs $(j,j+1)$ and $(k,k+1)$,
the corresponding four pair factors are distinct and
belong to (P-F57f3). The identity
\[
          y_{jk}+y_{j+1,k}+y_{j,k+1}=3\theta
                                      \pmod1
 \tag{P-F57f5}
\]
follows from $w_{j+1}=-2w_j$ and $w_{k+1}=-2w_k$.
The three displayed phases cannot all have distance
less than $2^{-u}$ from an integer. Thus at least one
of the four factors in each block is bad.
The blocks for distinct unordered pairs of index-pairs
are disjoint as sets of unordered factor pairs.
There are exactly $\binom h2$ such blocks. Therefore,
under this anchor condition, (P-F57f4) improves to
\[
       N_u\le\binom r2+3r+3
              -\binom{\lfloor(r-2)_+/2\rfloor}{2}.
 \tag{P-F57f6}
\]
The omitted factor $(2,3)$ is internal to one index-pair
and is not used in any of these four-cell blocks.

For any nonzero dyadic phase $y$ with denominator dividing
$M$, one has $\|y\|\ge1/M$ and
\[
 \log_2|1-e^{2\pi iy}|^{-1}
 \le-\log_2\|y\|
 \le\sum_{u=0}^{l-1}\mathbf1_{\{\|y\|<2^{-u}\}}.
 \tag{P-F57f7}
\]
The last inequality is the elementary layer-cake bound:
the number of integers $u$ with $0\le u< -\log_2\|y\|$
is its ceiling. The right side is nonnegative, so it
also bounds the sum over the increment's subset of active
pairs after adding back its one omitted factor.
All other increment factors number $2n-4$.
Each is nonzero modulo $M$, and hence has logarithmic
reciprocal at most $l$. Their total contribution is
less than $2n^2$.

Put $K=(\lfloor\log_2J\rfloor-2)_+$.
For every integer $1\le r\le K$ one has
$J>3\cdot2^r$, so the anchor condition of (P-F57f6)
holds at $u=l-r$.
Summing (P-F57f4)--(P-F57f7) gives
\[
 \log_2 I_n(a)
 \le\binom n3+\frac32l(l+1)+3l+2n^2
       -\sum_{r=1}^K
                  \binom{\lfloor(r-2)_+/2\rfloor}{2}.
 \tag{P-F57f8}
\]
The first four terms are at most $n^3/6+4n^2$ for $n\ge5$.
For $r\ge6$, the summand is at least $(r-5)^2/8$:
$\lfloor(r-2)/2\rfloor\ge(r-3)/2$, and the factors
in $h(h-1)/2$ are nonnegative in this range.
Consequently the sum is at least
$(K-5)_+^3/24\ge(\log_2J-8)_+^3/24$.

It remains to compare this upper bound with the selected
increment. Among its factors, the
$B=\binom{n-3}{2}-1$ active pair weights in (P-F57f3),
excluding $(2,3)$, are divisible by three.
At $a=b_n$ the factor for largest index $k$ is at most
$8\cdot2^k/M$, exactly as in (P-F56ah).
The sum of $n-1-k$ over these pairs is
$\binom{n-2}{3}-(n-4)$.
Every remaining increment factor has absolute value at
most two. There are $\binom{n-1}{2}$ factors in total.
It follows that
\[
 \begin{split}
 \log_2 I_n(b_n)
 &\ge\binom{n-2}{3}-(n-4)-2B-\binom{n-1}{2}\\
 &=n^3/6-3n^2+71n/6-11
 \ge n^3/6-3n^2.
 \end{split}
 \tag{P-F57f9}
\]
Subtracting proves (P-F57f1).
If $\log_2J\ge\delta n$ and $\delta n\ge16$,
then $(\log_2J-8)^3/24\ge\delta^3n^3/192$.
If also $n\ge2688/\delta^3$, then
$7n^2\le\delta^3n^3/384$.
These two inequalities prove (P-F57f2), with its explicit
sufficient dimension threshold.\hfill$\square$


\paragraph{Uniform contraction and the indices.}
For an odd residue $a\bmod M$, let $j$ be the least positive
absolute residue of $3a$ modulo $M$. Then $j$ is odd and
$1\le j\le M/2$, and $j=1$ holds exactly for
$a=b_n$ or $a=-b_n$.
Complex conjugation gives $J_n(a)=J_n(-a)$, so the
near estimate (P-F57e) covers every nonselected residue
with $j\le2^{\delta n}$.
For the others, the far-increment estimate just proved gives
a uniform bound $2^{-c_\delta n^3}$ for the same ratio,
with some constant $c_\delta>0$ and all sufficiently
large $n$.
It follows that there is a constant $c_*>0$ such that
\[
 \max_{\substack{a\bmod M,\ a\ {\rm odd}\\a\ne b_n,-b_n}}
          \frac{J_n(a)}{J_n(b_n)}
                         \le2^{-c_*n^2}
 \]
for all sufficiently large $n$.
Every parent class excluded from the two selected parent
classes in (P-F56ad) has two nonselected children.
Their normalized increments have sum at most
$2^{1-c_*n^2}\le1/2$ for all sufficiently large $n$.
Thus (P-F56ad) is proved with $\rho=1/2$.
The unconditional sibling bound and the recurrence
(P-F56ae)--(P-F56aj) yield $Q_n\to0$.
By (P-F56aa), this is exactly (P-F57a).

In particular $E_n\le |Z_n|/4$ eventually.
The now-discharged implication (P-F56n)--(P-F56p)
therefore gives the asserted exponential lower bounds
on both literal indices.
This last inference compares the isolated wave with actual
finite-dimensional symmetric powers. It does not infer an
index bound from the period alone or from an eventual
actual-support threshold.
\end{proof}
